\documentclass[11pt]{article}
\usepackage{fontspec}
\usepackage[french]{babel}
\usepackage{amsmath}
\usepackage{amssymb}
\usepackage{amsthm}
\usepackage{mathrsfs}
\usepackage{tikz-cd}
\usepackage{enumitem}

\newcommand{\printedpage}[1]{\marginpar{\footnotesize p.~#1}}
\newcommand{\msp}[1]{\marginpar{\footnotesize [#1]}}

\newcommand{\Ind}{\operatorname{Ind}}
\newcommand{\Filt}{\operatorname{Filt}}
\newcommand{\card}{\operatorname{card}}
\newcommand{\red}{\operatorname{\text{réd}}}
\newcommand{\Ob}{\operatorname{Ob}}
\newcommand{\Hom}{\operatorname{Hom}}
\newcommand{\uHom}{\underline{\operatorname{Hom}}}
\newcommand{\End}{\operatorname{End}}
\newcommand{\Ens}{\mathrm{Ens}}
\newcommand{\cM}{\mathcal{M}}
\newcommand{\cD}{\mathcal{D}}
\newcommand{\defeq}{\overset{\text{déf}}{=}}

\begin{document}

\printedpage{74}
% head: A. Grothendieck : Dérivateurs, Ch. XVIII, Édition des Universités de Montpellier II et Paris VII

\textbf{Lemme 4.8.10.} \textit{Soient $C$ une petite catégorie, $\pi_0$ un cardinal infini, $\pi_1 = \sup(\pi_0, \card\red C)$. Supposons $\Ind_{\pi_0}$ stable par $\varprojlim$ finies (resp. et d'un certain type). Alors pour tout $\pi \geqslant \pi_1$, $\Filt^{\pi}$ est stable par lesdites $\varprojlim$ . (Donc aussi $\cM_{\pi} = \Filt^{\pi}$, si $\pi = \pi^{\pi_0}$.) Plus} % [sic] \Ind_{\pi_0} stable -> \Ind_{\pi_0}C stable, the argument C appears to be omitted (cf. the following sentence)
\msp{page 209}
\textit{généralement, soit $J$ une petite catégorie telle que $\Ind_{\pi_0}C$ soit stable par $\varprojlim$ de type $J$, et soit $c = \card\red J$, supposons que $\pi^{c} = \pi$ (p.ex. $c$ fini). Alors $\Filt^{\pi}$ est stable dans $\cM$ par $\varprojlim$ de type $J$. (Donc il en est de même de $\cM_{\pi} = \Filt^{\pi}$, si de plus $\pi^{\pi_0} = \pi$.)}

On utilise le critère a$''$ de 4.7.1, savoir $\card F(x) \leqslant \pi$ pour tout $x \in C$, pour l'appartenance à $\Filt^{\pi}$. On peut supposer $J$ réduite, donc $\card\Ob J \leqslant c$. Comme le foncteur d'inclusion $\Ind_{\pi_0}(C) \longrightarrow C^{\wedge}$ commute aux $\varprojlim$ quelconques, le calcul de $\varprojlim\limits_{j\in J} x_j$ dans $C^{\wedge}$ fibre par fibre montre que dans $C^{\wedge}$, on a
\[
X \ \defeq\ \varinjlim_{j\in J} x_j \ = \ \underbrace{\prod_{j\in\Ob J} x_j}_{\text{produit dans } C^{\wedge}} \quad ,
\] % [sic] \varinjlim_{j\in J} x_j -> \varprojlim_{j\in J} x_j, the text concerns the projective limit (cf. the preceding sentence and (4.8.10.1))
d'où
\[
\card X(x) \ \leqslant\ \prod_{j\in\Ob J} \underbrace{\card x_j(x)}_{\leqslant\pi} \ \leqslant\ \pi^{c}\ ,
\]
d'où (sans autre hypothèse sur $\pi$ que $\pi \geqslant \pi_1$)
\[
\varprojlim_{j\in J} x_j \ \in\ \Ob\Filt^{\pi^{c}}\ , \tag*{(4.8.10.1)}
\]
d'où la conclusion voulue si $\pi^{c} = \pi$, q.e.d.

\msp{page 210}

Nous allons prouver les implications supplémentaires
\[
\text{b}'' \ \Longrightarrow\ \text{b}\ , \qquad \text{c} \ \Longrightarrow\ \text{b}\ , \qquad \text{c}'' \ \overset{?}{\Longrightarrow}\ \text{c}\ ,
\]
ce qui achève la démonstration.

L'implication b$''$ $\Longrightarrow$ a est contenue dans 4.6.3.

Pour l'implication c $\Longrightarrow$ b, on note que la stabilité de $\cM$ par petites $\varinjlim$ équivaut à l'assertion que toute petite limite projective de foncteurs $\cM \xrightarrow{\ \xi_i\ } \Ens$ représentables, est représentable. Or les objets de $\cM$ étant accessibles ($\cM$ étant accessible), les foncteurs représentables $\xi_i$ sont accessibles, donc aussi leur $\varprojlim$ (SGA 4 I 9.7). D'autre part, le foncteur $\xi$ est exact à gauche, comme $\varprojlim$ de foncteurs représentables donc exactes à gauche. Donc le fait que $\xi$ est représentable résulte du % [sic] exactes -> exacts, agreement with « foncteurs »

\textbf{Lemme 4.8.11.} \textit{Soit $\cM$ une catégorie accessible, stable par $\varprojlim$ finies, $\cM \longrightarrow$} Ens \textit{un foncteur. Pour que $\xi$ soit représentable, il faut et il suffit qu'il soit exact à gauche et accessible. Donc (SGA 4 I 8.12.8.1) si $\cM$ est stable par petites $\varprojlim$ , $\xi$ est représentable si et seulement s'il commute auxdites $\varprojlim$ , et qu'il soit accessible.}

% foot: Version 17 décembre 2025

\printedpage{75}
% head: A. Grothendieck : Dérivateurs, Ch. XVIII, Édition des Universités de Montpellier II et Paris VII

\msp{page 211}

Le ``il faut'' est remis pour mémoire, on vient de le voir comme conséquence du fait que tout objet de $\cM$ est accessible. Le ``il suffit'' est essentiellement SGA 4 I 8.13.2.b), sauf que dans loc. cit. on suppose $\cM$ munie d'une fc \textit{stricte.} Mais la démonstration donnée dans loc. cit. montre aussi pour une fcga, en choisissant dans loc. cit. (p. 137) $\pi$ tel que $\xi$ soit $\pi$-accessible, et de plus $\pi$ soit utile pour la filtration cardinale généralisée envisagée.

Reste à prouver c$''$ $\overset{?}{\Longrightarrow}$ c, sous réserve que ce soit bien vrai. Ce serait une conséquence du lemme suivant (où $\cD$, $\pi$ sont remplacés par $C$, $\pi_0$).

\textbf{Lemme 4.8.12} (?) (\footnote[88]{Lemme non prouvé ici -- cf. démonstration 4.11.7 (non achevée).}). \textit{Soient $C_0$ une petite catégorie, $\pi_0$ un cardinal infini. Si $C$ est stable par $\varprojlim$ finies, alors $\Ind_{\pi_0}C$ également.} % [sic] C_0 -> C, the statement then speaks only of C

L'assertion similaire pour $\Ind(C)$, au lieu de $\Ind_{\pi}C$, est établie dans SGA 4 I 8.9.1 c). Mais il n'est pas clair sans plus que $\Ind_{\pi}C$ soit stable dans $\Ind(C)$ par $\varprojlim$ finies, dans le cas envisagé. De façon précise, si $F = \varprojlim\limits_{j\in J} F_j$, les $F_j$ dans $\Ind_{\pi}C$,
\msp{page 212}
$F$ dans $\Ind(C) \subset C^{\wedge}$, on sait que $C/F$ est filtrant et il faut voir qu'il est filtrante grande devant $\pi_0$. On est ramené aux trois cas où $J = \emptyset$, $J = e \amalg e$, et $J = \vcenter{\hbox{\begin{tikzcd}[ampersand replacement=\&, column sep=small, row sep=small]
\bullet \arrow[dr] \& \& \bullet \arrow[dl] \\
\& \bullet \&
\end{tikzcd}}}$. Les deux premiers cas sont essentiellement triviaux (car $e$, et un produit de deux catégories grandes devant $\pi_0$, sont grandes devant $\pi_0$), le troisième cas par contre n'a pas l'air trivial a priori. C'est lié visiblement à la question des représentations indicielles de foncteurs $J \longrightarrow \Ind_{\pi_0}(C)$, pour $\card\red J \leqslant \pi_0$ -- i.e. la validité de 4.4.13 b) en dehors de l'hypothèse que $C$ soit stable par $\varinjlim$ de cardinal $\leqslant \pi_0$. Il nous faudra y revenir très prochainement, car cette question m'apparaît à présent une question-clef, pour les propriétés de stabilité de la notion de catégorie accessible (en dehors d'aucune hypothèse d'existence de $\varinjlim$ ou $\varprojlim$ ). % [sic] qu'il est filtrante -> qu'il est filtrant, agreement with « il » (C/F)

% marginal handwritten mark (curved stroke) in the left margin beside Proposition 4.8.13
\textbf{Proposition 4.8.13} (\footnote[89]{\textit{Pas prouvé.} Mais OK si $\pi_0 > \card\red C$, cf. 4.15.10, p. 462.}). \textit{Soit $C$ une catégorie équivalente à une petite catégorie, et soit $\pi_0 \geqslant \card\red C$ un cardinal infini. Alors $\cM = \Ind_{\pi_0}C = \cM_{\pi_0}$ (i.e. tous les objets de $\Ind_{\pi_0}C$ sont $\pi_0$-accessibles), et $\cM$ est une enveloppe de Karoubi de $C$.}

\msp{page 213}

Signalons tout de suite les corollaires suivants.

% marginal handwritten mark (curved stroke) in the left margin beside Corollaire 4.8.14
\textbf{Corollaire 4.8.14} (\footnote[90]{\textit{Pas prouvé.} Si finalement, cf. 4.15.10, et 4.15.11, pages 462-463.}). \textit{Soit $\cM$ une $\mathfrak{U}$-catégorie. Conditions équivalentes :}
\begin{enumerate}[leftmargin=*]
\item[\textit{a)}] \textit{$\cM$ est accessible, et est équivalent à une petite catégorie.}
\item[\textit{a$'$)}] \textit{$\cM$ est accessible, et il existe un cardinal infini $\pi_0$ tel que tous les objets de $\cM$ soient $\pi_0$-accessibles (i.e. $\cM = \cM_{\pi_0}$).}
\item[\textit{b)}] \textit{$\cM$ est équivalente à une petite catégorie, et est karoubienne, i.e. stable par facteurs directs} (\footnote[91]{\textbf{N.B.} Ceci implique que $\cM$ et $\cM^{\circ}$ sont accessible, et je conjecture que la réciproque est vraie également.}). % [sic] sont accessible -> sont accessibles, plural agreement (footnote 91)
\end{enumerate}

L'équivalence de a) et a$'$) est claire, puisque $\cM$ admet une filtration cardinale généralisée canonique $(\cM_{\pi})_{\pi\geqslant\pi_0}$, formée des catégories équivalentes à des petites catégories (condition
% foot: Version 17 décembre 2025
\printedpage{76}
% head: A. Grothendieck : Dérivateurs, Ch. XVIII, Édition des Universités de Montpellier II et Paris VII
a) des fcg), et telle que toute telle sous-catégorie soit contenue dans un des $\cM_{\pi}$ (résulte de la condition b) des fcg). Comme $\cM_{\pi_0}$ est équivalente à une petite catégorie, et est karoubienne (étant stable par $\varinjlim$ filtrantes grandes devant $\pi_0$), il est clair aussi que a$'$) $\Longrightarrow$ b). Reste à prouver b) $\Longrightarrow$ a$'$), or cela résulte aussitôt de 4.8.13 appliqué à $\cM$, qui montre de plus qu'on peut prendre pour $\pi_0$ n'importe quel cardinal infini tel que $\pi_0 \geqslant \card\red\cM$.

\msp{page 214}

\textbf{4.8.15.} \textsc{Démonstration de} 4.8.13. Nous allons prouver :

\textbf{Lemme 4.8.15.1.} \textit{Sous les conditions de 4.8.13, soit $F \in \Ob\cM$, $x \in \Ob C$, alors}
\[
\card F(x) \ \leqslant\ \pi_0\ . \tag*{$(*)$}
\]

Ceci signifie donc, en vertu de 4.7.7, a$''$ $\Longrightarrow$ b pour $\pi \geqslant \pi_1$, compte tenu qu'ici $\pi_J = \pi_0$, $\pi = \pi_0$, qu'on a $F \in \Filt^{\pi_0}$, d'où a fortiori $F \in \cM_{\pi_0}$, ce qui implique $\cM = \cM_{\pi_0}$, i.e. la proposition 4.8.13, compte tenu de 4.8.5. Donc tout revient à prouver $(*)$.

Écrivons
\[
F \ = \ \varinjlim_{i\in I} x_i\ , \qquad\quad \text{les } x_i \in \Ob C,\ I \text{ ensemble ordonné grand devant } \pi_0,
\]
d'où
\[
F(x) \ = \ \varinjlim_{i\in I} \Hom_C(x, x_i)\ .
\]
On peut supposer
\[
\card C \ = \ \pi_0\ ,
\]
quitte à remplacer $C$ par une catégorie réduite équivalente.

Pour tout $\xi \in \Ob C$, soit
\[
I_{\xi} \ = \ \{ i \in I \mid x_i = \xi \}\ ,
\]
alors
\[
I \ = \ \bigcup_{\xi\in\Ob C} I_{\xi}\ ,
\]
\msp{page 215}
réunion d'une famille de parties de $I$ de cardinal $\card\Ob C \leqslant \pi_0$. Comme $I$ est grand devant $\pi_0$, on voit que cela implique que l'ensemble des $I_{\xi}$ est cofinal dans $I$. Quitte à remplacer $I$ par $I_{\xi}$, on peut supposer $I_{\xi} = I$, i.e. que tous les $x_i$ sont égaux à un même élément $\xi$. Quitte à remplacer $C$ par la sous-catégorie pleine $C_0 = \{\xi\}$, on peut supposer $C$ réduite à un seul élément. (Car si on prouve que $F$, en tant qu'objet de $\Ind_{\pi_0}(C_0)$, est facteur direct d'un objet de $C_0$ (i.e. de $\xi$), il en sera de même de $F$ en tant qu'objet de $\Ind C$). Donc on peut supposer
\[
C \ = \ \{\xi\}\ ,
\]
donc $C$ donné par le monoïde
\[
M \ = \ \End_C(\xi)\ , \qquad\quad \card M \ \leqslant\ \pi\ ,
\]
et la donnée du système inductif $(x_i)_{i\in I}$ revient à la donnée d'un système d'éléments
\[
\alpha_{ji} \ \in\ M \qquad\quad (i, j \in I,\ i \leqslant j)
\]
% foot: Version 17 décembre 2025
\printedpage{77}
% head: A. Grothendieck : Dérivateurs, Ch. XVIII, Édition des Universités de Montpellier II et Paris VII
satisfaisant
\[
\alpha_{ii} \ = \ 1\ , \qquad\quad \alpha_{kj}\alpha_{ji} \ = \ \alpha_{ki} \ \ \text{si } i \leqslant j \leqslant k.
\]
On a
\[
F(\xi) \ = \ \varinjlim_{i\in I} M_i \qquad\quad
\left\{
\begin{array}{l}
\text{où } M_i = M \ \forall\, i \in I, \text{ et } M_i \longrightarrow M_j \\
\text{est donné par} \\[2ex]
\qquad\qquad u \longmapsto \alpha_{ji}u\ . \\[4ex]
\end{array}
\right.
\]

\msp{page 216}

J'ai pu me ramener, par des arguments similaires (\footnote[92]{Pour ces réductions, cf. 4.15.8 (p. 453-463). L'ultime réduction se trouve dans 4.15.8.9 (p. 461), et son corollaire 4.15.10.}), au cas où le monoïde $M = \End(\xi)$ est un monoïde de monomorphismes, i.e.
\[
uv \ = \ uv' \ \ \Longrightarrow\ \ v \ = \ v' \qquad\quad \text{si } u, v, v' \in \cM\ ,
\] % [sic] u, v, v' \in \cM -> u, v, v' \in M, the elements belong to the monoid M
ce qui implique que les flèches de transition dans le système inductif des $M_i = M$ sont des monomorphismes. Il s'ensuit que l'on ne peut avoir $\card F(\xi) \leqslant \pi_0$ ($I$ étant grand devant $\pi_0$) que si $\alpha_{ji_0}$ est un isomorphisme pour $j \geqslant i_0$, et $i_0$ assez grand, i.e. si le système inductif induit pour $j \geqslant i_0$ est constant, a fortiori que le système inductif envisagé est essentiellement constant. Mais je viens de passer près de quatre heures pour essayer de prouver ce point, sans y être arrivé ! Cela semble plus délicat à prouver que je ne le prévoyais, mais je crois toujours que c'est bel et bien vrai.

\msp{page 217}

\subsection*{4.9\quad Le théorème de représentation indicielle pour $\Ind_{\pi}(C)$}

\textbf{4.9.1} (\footnote[93]{\texttt{[Les numéros 4.9.1 à 4.9.6 sont utilisés deux fois, 4.9.3 trois fois.]}}). Soient $C$ une catégorie équivalente à une petite catégorie, $\pi_0$ un cardinal infini,
\[
\cM \ = \ \Ind_{\pi_0}(C) \ \subset\ C^{\wedge}\ , \tag*{(4.9.1)}
\]
$J$ une catégorie équivalente à une petite catégorie. Considérons l'inclusion pleinement fidèle $C \hookrightarrow \cM$, et l'inclusion correspondante de $\mathfrak{U}$-catégories
\[
\uHom(J, C) \ \hookrightarrow\ \uHom(J, \cM) \qquad\quad \text{(pleinement fidèle).} \tag*{(4.9.2)}
\]
On sait, comme $C \subset \cM_{\pi_0}$, que les objets de $\uHom(J, C)$ donnent des objets $\pi_0$-\textit{accessibles} de $\uHom(J, \cM)$, du moins si on suppose
\[
\card\red J \ \leqslant\ \pi_0\ , \tag*{(4.9.3)}
\]
ce que nous ferons désormais. (Cf. 4.5.8, 2°, p. 119.) Comme le deuxième membre de (4.9.2) satisfait $\mathrm{L}_{\pi_0}$ ($\cM$ y satisfait), il s'ensuit que ce foncteur (4.9.2) se prolonge en
\[
\Ind_{\pi_0}(\uHom(J, C)) \ \hookrightarrow\ \uHom(J,\ \underbrace{\cM}_{=\,\Ind_{\pi_0}(C)}\ )\ , \tag*{(4.9.3)}
\]
% foot: Version 17 décembre 2025
\printedpage{78}
% head: A. Grothendieck : Dérivateurs, Ch. XVIII, Édition des Universités de Montpellier II et Paris VII
et ce foncteur est \textit{pleinement fidèle} grâce au résultat d'accessibilité qu'on vient d'énoncer. Je m'intéresse à établir des cas où ce foncteur est essentiellement surjectif, i.e. est une \textit{équivalence} de catégories. Notons que pour que
\[
\varphi \ :\ J \longrightarrow \cM
\]
appartienne à l'image essentielle, il faut et il suffit qu'il existe un foncteur
\[
\Phi \ :\ I \ \longrightarrow\ \uHom(J, C)\ ,\ I \text{ grande devant } \pi_0 \text{ (et on peut supposer } I \text{ ordonnée)}\ , \tag*{(4.9.4)}
\]

\msp{page 218}

tel qu'on ait
\[
\varphi \ \xleftarrow{\ \sim\ }\ \varinjlim_{i\in I} \Phi(i)\ , \tag*{(4.9.5)}
\]
ou ce qui revient au même, tel qu'on ait
\[
\varphi(j) \ \xleftarrow{\ \sim\ }\ \varinjlim_{i\in I} \Phi(i, j)\ , \qquad\quad \text{pour } j \in J\ , \tag*{(4.9.6)}
\]
où on a posé
\[
\Phi(i, j) \ = \ \Phi(i)(j)\ ,
\]
interprétant donc $\Phi$ comme un foncteur
\[
\Phi \ :\ I \times J \ \longrightarrow\ C\ ; \tag*{(4.9.7)}
\]
la loi fonctorielle de $\varphi(j)$ donnée par (4.9.6) étant évidente. En d'autres termes, les $\varphi$ dans l'image essentielle de (4.4.3) sont ceux qui admettent une \textit{représentation indicielle} par un foncteur $\Phi$ sous forme (4.9.4), ou sous forme (4.9.7), avec \textit{de plus} $I$ grand devant $\pi_0$. Je rappelle a ce sujet la % [sic] a ce sujet -> à ce sujet, missing accent

\textbf{Proposition 4.9.2.} \textit{Sous les conditions précédentes sur $C$, $\pi_0$, $J$ (notamment} $\card\red J \leqslant \pi_0$\textit{), si $\varphi : J \longrightarrow \cM \defeq \Ind_{\pi_0}(C)$, pour que $\varphi$ appartienne à l'image essentielle du foncteur pleinement fidèle (4.9.3), il faut et il suffit qu'elle satisfasse les deux conditions suivantes :}

\msp{page 219}

\begin{enumerate}[leftmargin=*]
\item[\textit{a)}] $\uHom(J, C)/\varphi$ \textit{est grande devant $\pi_0$.}
\item[\textit{b)}] \textit{Pour tout $j \in J$, le foncteur}
\[
\uHom(J, C)/\varphi \ \longrightarrow\ C/\varphi(j) \qquad\quad
\left\{
\parbox{0.35\textwidth}{\textit{foncteur entre catégories filtrantes grandes devant $\pi_0$, en vertu de a),}}
\right. \tag*{(4.9.2.1)}
\]
\textit{induit par le foncteur}
\[
\left\{
\begin{array}{rcl}
\uHom(J, C) & \longrightarrow & C \\
F & \longmapsto & F(j)\ ,
\end{array}
\right. \tag*{(4.9.2.2)}
\]
\textit{est} cofinal.
\end{enumerate}
% foot: Version 17 décembre 2025

\end{document}
